\section{Substitutions on compact alphabets} \label{sec:alph}


In this section, we adapt the notations and definitions in \cite{MRW2}. 
\subsection{Subshifts and substitutions}

%

Let $\mathcal{A}$ be a compact Hausdorff space. The set $\A$ will be our \emph{alphabet}. Its elements are called \emph{letters}.  Let $\mathcal{A}^n$ be the collection of all words of length $n$ over $\mathcal{A}$, i.e., $w=w_1w_2\cdots w_n$, with $w_i\in\mathcal{A}$. 
We say that a %
word $u$ is a subword of $w$ (which we denote by $u\triangleleft w$) if there exist $1\leqslant i \leqslant j\leqslant n$ such that $u=w_iw_{i+1}\cdots w_{j}= w_{[i,j]}$. Note that the relation of ``being a subword'' can be easily generalized when $w$ is one-sided infinite or bi-infinite.
The set $\mathcal{A}^{+}=\bigcup_{n\geqslant 1} \mathcal{A}^n$ 
of all (non-empty) finite words over $\mathcal{A}$ is also a topological space (since the sets $\mathcal{A}^n$ are naturally endowed with the product topology), where the topology on $\mathcal{A}^{+}$ is that of a disjoint union. Let $\mathcal{A}^{\ast}=\mathcal{A}^{+}\cup \left\{\varepsilon\right\}$, where $\varepsilon$ is the empty word. 
$\A^*$ is a free monoid with concatenation as the binary operation. 

\begin{definition}
A {\it substitution} $\varrho$ is a continuous map from $\mathcal{A}$ to $\mathcal{A}^{+}$. 
\end{definition}

For every letter $a\in \mathcal A$, we call $\varrho^{k}(a)$ the \emph{level-$k$ supertile} of type $a$. 
We define the \emph{language} $\mathcal{L}(\varrho)$ of the substitution to be the set 

%
%
%
%
%
\[
\mathcal{L}(\mathcal{\varrho})=\overline{\left\{
u\in \mathcal{A}^{\ast}\mid u \triangleleft \varrho^{k}(a), k\in\mathbb{N}_0,a\in\mathcal{A}\right\}}. 
\]
So, the language $\mathcal{L}(\varrho)$ is the closure of the family of all subwords of supertiles.

This allows one to define a subshift associated to $\varrho$. Let $\mathcal{A}^{\mathbb{Z}}$ be the set of all bi-infinite sequences over $\mathcal{A}$, also called the \emph{full shift} over $\mathcal{A}$, which is compact with respect to the Tychonoff topology. Note that the left shift map $\sigma\colon \mathcal{A}^{\mathbb{Z}}\to \mathcal{A}^{\mathbb{Z}}$, defined pointwise via $\sigma(x)_n=x_{n+1}$, is a homeomorphism. 
A nonempty subset $X\subseteq \mathcal{A}^{\mathbb{Z}}$ is called a \emph{subshift} if it is closed and shift-invariant, i.e., $\sigma(X)\subseteq X$. 
We now define the subshift associated to $\varrho$ via its corresponding language, i.e., $X_{\varrho}:=\left\{x\in\mathcal{A}^{\mathbb{Z}}\mid u\in\mathcal{L}(\varrho), \text{ for all }u\triangleleft x\right\}$; see \cite{MRW2}. It can be shown that $\sigma(X_{\varrho})=X_{\varrho}$. 
A substitution $\varrho$ is said to be \emph{recognizable} if for every power $k\geqslant 1$, every element $x$ of the subshift admits a unique decomposition into level-$k$ supertiles of $\varrho$.

\begin{definition}
A substitution is called \emph{primitive} if for every non-empty open set $U\subset \mathcal{A}$, there exists $k(U)\in \mathbb{N}$ such that $\varrho^{k(U)}(a)$ contains an element in $U$ for all $a\in\mathcal{A}$. 
\end{definition}

Primitivity has several immediate consequences for the corresponding subshift, some of which we mention in the following result, see \cite[Thm. 3.30]{MRW2}. %

\begin{prop}
Let $\varrho$ be a primitive substitution over a compact Hausdorff alphabet. Then 
\begin{enumerate}
\item[\textnormal{(i)}] The subshift $X_{\varrho}$ is non-empty.
\item[\textnormal{(ii)}] The topological dynamical system $(X_{\varrho},\sigma)$ is minimal, i.e., every element $x\in X_{\varrho}$ has a dense orbit. 
\end{enumerate}
\end{prop}

\subsection{The substitution operator}

Let $\varrho\colon \mathcal{A}\to \mathcal{A}^{+}$ be a substitution over a compact Hausdorff alphabet. One can associate to it a substitution operator $M:=M_{\varrho}$ on the space $C(\mathcal{A})$ of real-valued continuous functions over $\mathcal{A}$
via
\[
Mf(a)=\sum_{b\,\triangleleft\,\varrho(a)} f(b),
\]
where the summation goes over all letters $b$ and $\varrho(a)$ is seen as a multiset, i.e. the summation takes into account multiplicities of letters in $\varrho(a)$ as well. 
This is the compact alphabet analogue of the transpose of the substitution matrix in the finite alphabet setting. This is a bounded operator because the length of $\varrho(a)$ is bounded. 
Let $\mathcal{K}$ be the positive cone of $C(\mathcal{A})$, i.e., all functions with $f(a)\geqslant 0$ for all $a\in\mathcal{A}$. One has $M\mathcal {K}\subseteq \mathcal{K}$, that is, $M$ is a positive operator on $C(\mathcal{A})$. 

The generalization of the Perron--Frobenius eigenvalue in this setting is the spectral radius $r(M)$ of $M$. Since $M$ is positive, $r(M)$ is always in the spectrum of $M$ but it is not guaranteed that it is an eigenvalue.
We refer the reader to \cite{Sch} for a background on positive operators on Banach lattices. 
 For the substitution operator, one has the following bounds for $r(M)$. 

\begin{lem}[\cite{MRW2}, Lemma 4.23]\label{lem:spec-rad}
For all $k\in\mathbb{N}$ one has
\[
\min_{a\in\mathcal{A}} |\varrho^{k}(a)|\leqslant r(M)^{k}\leqslant \max_{a\in\mathcal{A}} |\varrho^{k}(a)|. 
\]
\end{lem}

An operator $T$ with $r(T)=1$ over a Banach space  is called \emph{quasicompact} if there exists $k\in\mathbb{N}$ and a compact operator $K$ such that $\|T^k-K\|<1$, where $\| \cdot \|$ is the operator norm induced by the sup-norm on $C(\mathcal{A})$.

\begin{prop}[\cite{MRW2}, Theorem 6.8(1)] \label{prop:quasicompact}
Let $\varrho$ be a primitive substitution on a compact Hausdorff alphabet $\mathcal{A}$ for which there exists $k\in\mathbb{N}$ and a finite subset $F\subset \mathcal{A}$ of isolated points such that
\[
\left|\left\{b\in\varrho^{k}(a)\mid b\notin F\right\}\right|<r(M)^k
\]
for all $a\in\mathcal{A}$. Then the normalized substitution operator $T:=M/r(M)$ is quasicompact. 
\end{prop}

\begin{thm}[\cite{MRW2}, Theorem 6.7(3,4,5)]\label{thm:unique-ergod}
Let $\varrho$ be a primitive substitution on a compact Hausdorff alphabet. Suppose further that $T=M/r(M)$ is quasicompact. Then the following hold
\begin{enumerate}[label={{\rm (\roman*)}}]
\item $\varrho$ admits a unique natural tile length function which is strictly positive.
\item The corresponding tiling dynamical system $(\Omega_{\varrho},\mathbb{R})$ is uniquely ergodic.
\item The corresponding symbolic dynamical system $(X_{\varrho},\sigma)$ is  uniquely ergodic. 
\end{enumerate}
\end{thm}
